\chapter{Tìm hiểu về boson Higgs}

\lectureoverview{Cơ chế Higgs dễ hiểu nhất khi ta trình bày các thành phần theo đúng trình tự chúng trở nên cần thiết. Trước hết, ta xét chuyển động và trường đã được lượng tử hóa, rồi biến một cái bát thông thường thành thế năng hình mũ Mexico; từ trạng thái ngưng tụ hình thành, ta hiểu được vì sao fermion và boson yếu có thể có khối lượng. Chỉ sau đó ta mới xác định boson Higgs vật lý là dao động theo phương bán kính của trạng thái ngưng tụ ấy, rồi lần theo các quá trình tạo sinh và phân rã đã giúp khám phá ra nó vào năm 2012.}

\section{Cơ chế đằng sau những dòng tít}

Việc phát hiện một hạt mới có khối lượng gần $125\,\mathrm{GeV}$ quả thật khiến người ta phấn khích, nhưng phấn khích thôi thì chưa phải là hiểu biết. Câu hỏi hữu ích không phải là Higgs có quan trọng hay không, cũng không phải liệu ta có thể dùng lời lẽ hoa mỹ để gắn nó với nguồn gốc vũ trụ hay không. Câu hỏi hữu ích hẹp hơn nhiều: hiện tượng này vận hành cụ thể ra sao, và vì sao Mô hình Chuẩn cần đến nó?

Tên gọi này cũng cần được nói rõ đôi chút. Peter Higgs là một trong số những người đã nhận ra vật lý đằng sau cơ chế, còn Fran\c{c}ois Englert và Robert Brout đã độc lập đi đến ý tưởng cốt lõi. Vì vậy, cụm từ \emph{cơ chế Brout----Englert} phản ánh lịch sử chính xác hơn, dù phần dưới đây sẽ dùng tên gọi ngắn hơn.

Lời giải thích gồm nhiều phần. Trước tiên, ta cần hai khái niệm cơ bản của cơ học lượng tử và một ý tưởng của lý thuyết trường cổ điển. Tiếp đó, ta cần một chân không mà trạng thái năng lượng thấp nhất chứa một trường khác không. Sau đó, ta có thể điểm qua các hạt trong Mô hình Chuẩn và thấy vì sao phiên bản sơ khai của lý thuyết khiến những hạt liên quan không có khối lượng. Boson Higgs chỉ xuất hiện gần cuối. Nó không phải toàn bộ cơ chế, mà là một kiểu dao động cụ thể còn lại sau khi cơ chế đã vận hành.

\section{Hai sự thật lượng tử và một trường}

Sự thật lượng tử thứ nhất là sự lượng tử hóa. Mômen động lượng là ví dụ điển hình dễ thấy vì nó chỉ nhận những giá trị rời rạc,
\begin{equation}
L=n\hbar,
\qquad
n=0,\pm1,\pm2,\ldots .
\end{equation}
Mômen động lượng bán nguyên cũng tồn tại, nhưng lập luận trước mắt không phụ thuộc vào chúng. Điều cốt yếu là chuyển động cổ điển biến thiên liên tục trở thành một dãy các trạng thái lượng tử.

Sự thật thứ hai là nguyên lý bất định. Ta sẽ để dành nó đến khi xuất hiện một đường tròn liên tục gồm các chân không khả dĩ. Khi đó, nguyên lý này sẽ ngăn ta xem chân không như một quả cầu cổ điển nhỏ nằm yên tại một điểm được xác định chính xác.

Giữa hai sự thật lượng tử này là khái niệm trường của vật lý cổ điển. Trường gán một đại lượng cho mỗi điểm trong không gian và thời gian. Trường điện, từ và hấp dẫn là những ví dụ quen thuộc. Một vùng có thể trông như trống rỗng mà vẫn chứa trường. Hãy hình dung hai bản tích điện ở xa nhau đến mức không thể nhìn thấy bản nào. Dẫu vậy, điện trường gần như đồng đều giữa chúng vẫn có thực tại chỗ. Vì thế, ``không gian trống rỗng'' không nhất thiết có nghĩa là ``mọi trường đều bằng không.'' Điều đó có nghĩa là các trường ở trạng thái có năng lượng thấp nhất trong số những trạng thái sẵn có.

Giả sử năng lượng của một trường nhỏ nhất tại một giá trị nào đó và tăng lên khi trường lệch khỏi giá trị ấy. Nếu trường bị nhiễu động, nó dao động quanh điểm cực tiểu. Cơ học lượng tử lượng tử hóa các dao động đó, và những lượng tử riêng lẻ của chúng là các hạt. Ta sẽ dùng xuyên suốt bài giảng cách chuyển đổi giữa hai cách mô tả của lý thuyết trường sau đây:
\begin{equation}
\text{particle}
\quad\longleftrightarrow\quad
\text{quantum of a field vibration}.
\end{equation}

\section{Một cái bát trong không gian trường}

Với một thành phần trường thực $\phi$, cảnh quan năng lượng đơn giản nhất có dạng một cái bát,
\begin{equation}
V(\phi)=\frac{1}{2}m^{2}\phi^{2}.
\end{equation}
Đáy bát nằm tại $\phi=0$. Những dao động nhỏ quanh đáy có một tần số xác định, và các lượng tử của chúng ứng xử như những hạt có khối lượng $m$.

Giờ thay trường đơn bằng hai thành phần, $(\phi_1,\phi_2)$. Năng lượng có thể chỉ phụ thuộc vào bán kính trong mặt phẳng nội tại này,
\begin{equation}
r^{2}=\phi_1^{2}+\phi_2^{2},
\qquad
V=V(r).
\end{equation}
Đây là một cái bát trong \emph{không gian trường}, chứ không phải một chỗ trũng trong không gian thông thường. Phép quay cặp $(\phi_1,\phi_2)$ không làm thay đổi năng lượng. Chiếc mũ dùng trong bài giảng giúp ta hình dung cùng tính đối xứng ấy: xoay nó quanh trục thì hình dạng vẫn không đổi.

\begin{figure}[tbp]
\centering
\includegraphics[width=0.95\textwidth]{lecture_144_figure_02.png}
\caption{Cái bát đối xứng thông thường và chiếc mũ vật lý dùng để giúp hình dung tính đối xứng quay trong không gian trường.}
\lectureframe{00:10:02}
\end{figure}

Phép tương tự không chỉ còn là hình học khi trường chuyển động. Hãy làm trường lệch khỏi tâm rồi cho nó chuyển động tuần hoàn quanh đáy bát. Chuyển động tuần hoàn này mang mômen động lượng trong không gian trường. Đây không phải mômen động lượng của một vật quay trong phòng thí nghiệm, mà là mômen động lượng nội tại. Khi được lượng tử hóa, nó trở thành mô hình cho một điện tích được bảo toàn,
\begin{equation}
Q=nq_0,
\qquad
n\in\mathbb{Z}.
\end{equation}
Phương trình này chỉ có tính sơ đồ. Ý tưởng vật lý cốt yếu là có thể hình dung một kích thích của trường mang điện như chuyển động quay trong mặt phẳng nội tại. Điện tích dương và âm tương ứng với hai chiều quay ngược nhau; không quay tương ứng với điện tích bằng không.

\section{Lật úp chiếc mũ}

Cái bát thông thường chỉ có một điểm thấp nhất. Vật lý Higgs bắt đầu khi lật úp chiếc mũ. Tâm lúc này không ổn định, còn cả một rãnh tròn có cùng năng lượng cực tiểu. Một cách hiện thực hóa toán học tiêu chuẩn là
\begin{equation}
V(\phi_1,\phi_2)
=
\lambda
\left(
  \phi_1^{2}+\phi_2^{2}-v^{2}
\right)^{2},
\qquad \lambda>0,
\end{equation}
trong đó các điểm cực tiểu thỏa mãn
\begin{equation}
\phi_1^{2}+\phi_2^{2}=v^{2}.
\end{equation}
\editorialnote{Thế năng bậc bốn tường minh là cách hoàn chỉnh giải tích tiêu chuẩn cho hình học mũ Mexico được vẽ trong bài giảng; biểu thức này không được chép lại từ công thức trên bảng.}

\begin{figure}[tbp]
\centering
\includegraphics[width=0.95\textwidth]{lecture_144_figure_03.png}
\caption{Đỉnh ở tâm và rãnh tròn của thế năng hình mũ lật ngược. Bản phác thảo đầy đủ trên bảng thay cho một hình cắt trước đó đã bỏ sót phần lớn hình dạng.}
\lectureframe{00:13:10}
\end{figure}

Chân không giờ đây không còn có trường bằng không. Nó phải nằm ở đâu đó trên rãnh,
\begin{equation}
(\phi_1,\phi_2)
=
v(\cos\theta,\sin\theta).
\end{equation}
Các phương trình vẫn giữ tính đối xứng quay, nhưng mỗi chân không cụ thể chọn ra một góc $\theta$. Đây là sự phá vỡ đối xứng tự phát: định luật đối xứng, còn trạng thái năng lượng thấp nhất được chọn thì không.

Giờ đây có hai kiểu chuyển động khác nhau về bản chất. Chuyển động theo phương bán kính vào trong hoặc ra ngoài sẽ đi lên thế năng và tốn năng lượng. Chuyển động quanh rãnh chỉ làm thay đổi $\theta$ và, ở mức cổ điển, không tốn năng lượng thế. Vì vậy, có thể gắn một chuyển động góc giống như chuyển động của điện tích rất lớn với mức tăng năng lượng thế rất nhỏ.

\begin{figure}[tbp]
\centering
\includegraphics[width=0.95\textwidth]{lecture_144_figure_04.png}
\caption{Một điểm được đánh dấu trên rãnh biểu thị một chân không cụ thể được chọn từ đường tròn liên tục gồm các điểm cực tiểu có cùng năng lượng.}
\lectureframe{00:16:14}
\end{figure}

\section{Nguyên lý bất định biến chân không thành trạng thái ngưng tụ}

Theo vật lý cổ điển, ta có thể đặt trường tại một điểm trên vành, cho vận tốc góc của nó bằng không rồi xem như vấn đề đã xong. Cơ học lượng tử không cho phép kết luận chắc chắn đến vậy. Vị trí và động lượng tuân theo
\begin{equation}
\Delta x\,\Delta p\gtrsim\frac{\hbar}{2}.
\end{equation}
Phát biểu tương tự trên rãnh tròn là: định xứ góc của trường khiến mômen động lượng nội tại liên hợp với nó trở nên bất định. Theo cách trình bày sơ đồ của bài giảng,
\begin{equation}
\Delta\theta\,\Delta L_{\mathrm{int}}
\gtrsim
\frac{\hbar}{2}.
\end{equation}
\editorialnote{Toán tử góc có những điểm phức tạp về mặt kỹ thuật, vì vậy hệ thức này biểu đạt nguyên lý bất định trong bài giảng theo dạng bán cổ điển cục bộ, chứ không phải một đẳng thức toán tử chính xác trên toàn miền.}

Vì thế, chân không tập trung gần một điểm trên rãnh không phải là trạng thái có điện tích nội tại xác định. Nó chứa một chồng chập rộng các trạng thái riêng của điện tích. Bài giảng làm rõ điều này bằng một phép thử thực tế: thêm một đơn vị điện tích hoặc bớt đi một đơn vị hầu như không làm thay đổi trạng thái. Có thể hình dung sơ đồ như sau
\begin{align}
\ket{\Omega}
&\sim
\sum_n c_n\ket{n},
\\
c_{n+1}
&\simeq
c_n
\quad\text{within the occupied range}.
\end{align}
Khi đó, dịch $n$ đi một đơn vị chủ yếu chỉ đánh lại nhãn cho những số hạng vốn đã có mặt. Đây là ý nghĩa của việc chân không là một \emph{trạng thái ngưng tụ}: nó chứa một số lượng không xác định các lượng tử liên quan, nên thêm hoặc bớt một lượng tử cũng không tạo ra khác biệt vĩ mô có thể quan sát được.

Chân không điện yếu thực sự không phải là trạng thái ngưng tụ của điện tích thông thường. Nếu là như vậy, tính chất điện từ của không gian trống rỗng sẽ khác hẳn. Tuy nhiên, chất siêu dẫn là một ví dụ cụ thể về trạng thái ngưng tụ mang điện. Vì vậy, nó tạo nên cầu nối vật lý hữu ích giữa hình ảnh trừu tượng về chiếc mũ Mexico và nhận định rằng chân không có thể mang cấu trúc lượng tử kết hợp.

Như vậy, phần thứ nhất đã hoàn tất. Ta đã biết trạng thái năng lượng thấp nhất không nhất thiết phải có trường bằng không, và bất định lượng tử có thể khiến trạng thái ấy không nhạy với việc thêm hoặc bớt từng lượng tử riêng lẻ. Giờ ta hỏi những hạt nào cần có phông nền như vậy.

\section{Mô hình Chuẩn sơ khai}

Giới hạn trên tự nhiên để so sánh khối lượng các hạt cơ bản là khối lượng Planck. Khi năng lượng tập trung đủ cao, ta không thể bỏ qua hấp dẫn; đến một mức nào đó, vật thể phù hợp hơn với mô tả lỗ đen thay vì hạt cơ bản thông thường. Giá trị số là
\begin{equation}
m_{\mathrm P}
=
\sqrt{\frac{\hbar c}{G}}
\approx
2.2\times10^{-8}\,\mathrm{kg},
\end{equation}
xấp xỉ khối lượng của một hạt bụi nhỏ. Các hạt cơ bản đã biết nhẹ hơn con số này vô cùng nhiều. Theo ước lượng về thang khối lượng dùng trong bài giảng, khối lượng của chúng chỉ đạt khoảng $10^{-17}$ khối lượng Planck.

Nhận xét đó không chứng minh rằng không tồn tại hạt cơ bản nặng hơn. Máy gia tốc có xu hướng phát hiện những hạt đủ nhẹ để tạo ra với năng lượng sẵn có. Vì thế, phổ các hạt đã biết một phần được quyết định bởi tầm với của những cỗ máy dùng để tìm kiếm chúng. Có thể còn nhiều hạt nặng hơn nằm ngoài khả năng tiếp cận hiện nay.

Các hạt quan sát được chia thành fermion và boson. Fermion gồm electron, muon, tau, các neutrino của chúng và sáu hương quark. Boson gồm photon, gluon và các boson yếu $W$ và $Z$. Gluon liên kết các quark thành proton và neutron, nên giữ vai trò thiết yếu trong cấu trúc của các hạt nhân thông thường.

\begin{figure}[tbp]
\centering
\includegraphics[width=0.95\textwidth]{lecture_144_figure_05.png}
\caption{Bảng kiểm kê đầu tiên về Mô hình Chuẩn trong bài giảng, với dải khối lượng khổng lồ được vẽ phía trên các cột fermion và boson.}
\lectureframe{00:24:22}
\end{figure}

Lý thuyết không chỉ là một danh sách hạt mà còn là tập hợp các tương tác được phép. Electron có thể phát ra photon,
\begin{equation}
e\longrightarrow e+\gamma.
\end{equation}
Quark mang cả điện tích điện lẫn điện tích màu, nên có thể phát ra photon hoặc gluon,
\begin{equation}
q\longrightarrow q+\gamma,
\qquad
q\longrightarrow q+g.
\end{equation}
Electron không thể phát ra gluon, còn neutrino không phát ra photon cũng như gluon. Các fermion có số lượng tử yếu thích hợp có thể tương tác thông qua boson $Z$,
\begin{equation}
f\longrightarrow f+Z.
\end{equation}
Các boson $W$ và những phản ứng dòng mang điện cũng thuộc Mô hình Chuẩn, nhưng bài giảng lược bỏ chúng để chỉ cần dùng một boson yếu trung hòa làm trọng tâm lập luận.

\begin{figure}[tbp]
\centering
\includegraphics[width=0.95\textwidth]{lecture_144_figure_06.png}
\caption{Các đỉnh phát xạ photon, gluon và $Z$ nằm cạnh từ ``không khối lượng.'' Các quá trình này tóm tắt lý thuyết sơ khai trước khi đưa trạng thái ngưng tụ vào.}
\lectureframe{00:29:20}
\end{figure}

Photon không có khối lượng. Trong một lý thuyết chuẩn không phá vỡ có cùng dạng tổng quát, gluon và các boson chuẩn yếu cũng không có khối lượng. Tương tự, lý thuyết điện yếu chưa biến đổi không thể đơn giản thêm khối lượng fermion bằng tay mà vẫn giữ tính đối xứng chuẩn. Còn thiếu một thành phần. Trước khi bổ sung nó, ta cần xem trong một bối cảnh đơn giản hơn trường nền có thể làm thay đổi khối lượng hạt ra sao mà không tác dụng như lực ma sát.

\section{Trường nền có thể làm thay đổi khối lượng ra sao}

Hãy xem phân tử nước như một lưỡng cực điện: một đầu hơi dương, đầu kia hơi âm. Trong không gian trống, hai hướng định hướng ngược nhau có năng lượng tương đương. Bài giảng vui đùa đặt cho hai hướng những tên gọi khác nhau, ``nước'' và ``rượu whisky,'' chỉ để có thể theo dõi chúng như những trạng thái nội tại riêng biệt.

Đặt phân tử giữa hai bản tụ điện. Nếu mômen lưỡng cực của nó là $\bm d$ và điện trường là $\bm E$, năng lượng tương tác bằng
\begin{equation}
U=-\bm d\mathbin{\cdot}\bm E.
\end{equation}
Với hai hướng định hướng thẳng hàng, năng lượng trở thành
\begin{equation}
U_{\parallel}=-dE,
\qquad
U_{\mathrm{anti}}=+dE.
\end{equation}
Hệ thức Einstein $E=mc^2$ biến độ chênh lệch năng lượng thành độ chênh lệch khối lượng,
\begin{equation}
m_{\mathrm{anti}}-m_{\parallel}
=
\frac{2dE}{c^{2}}.
\end{equation}
Không có gì cọ xát vào phân tử. Chỉ đơn giản là một trạng thái nội tại có năng lượng, và do đó có khối lượng, lớn hơn trạng thái kia khi trường nền hiện diện.

\begin{figure}[tbp]
\centering
\includegraphics[width=0.95\textwidth]{lecture_144_figure_07.png}
\caption{Hai hướng định hướng ngược nhau của phân tử trong điện trường giữa hai bản tụ. Trường làm tách mức năng lượng của chúng mà không tạo ra lực cản.}
\lectureframe{00:34:35}
\end{figure}

Điều này bác bỏ hai phép so sánh dễ gây hiểu lầm. Trường Higgs không phải là mật đường mà các hạt phải ráng sức xuyên qua; hạt có khối lượng cũng không phải xe ủi tuyết bị vật chất tích tụ phía trước làm chậm lại. Trong một trường đồng đều theo không gian, tổng lực tác dụng có thể bằng không. Hiệu ứng ở đây là sự dịch chuyển năng lượng nội tại, chứ không phải ma sát hay mất động lượng.

Cũng có thể kể lại ví dụ này bằng ngôn ngữ lượng tử. Trường điện cổ điển là một trạng thái kết hợp, có thể được mô tả trong phạm vi đang xét như một trạng thái ngưng tụ photon. Phân tử có thể phát ra một photon vào nền đó hoặc hấp thụ một photon từ nền. Vì thêm hoặc bớt một lượng tử riêng lẻ hầu như không làm thay đổi trạng thái kết hợp có số chiếm dụng lớn, đường biểu diễn photon có thể kết thúc trong trạng thái ngưng tụ. Dấu gạch chéo vẽ trên bảng biểu thị photon biến mất vào trường nền, chứ không phải sự vi phạm định luật bảo toàn năng lượng.

\begin{figure}[tbp]
\centering
\includegraphics[width=0.95\textwidth]{lecture_144_figure_08.png}
\caption{Điện trường giữa hai bản tụ được vẽ lại dưới dạng trạng thái ngưng tụ photon. Photon do lưỡng cực phát ra bị hấp thụ vào trường nền vốn không nhạy với việc thêm một lượng tử.}
\lectureframe{00:38:45}
\end{figure}

Đây là khuôn mẫu ta cần. Nếu trạng thái ngưng tụ có thể hấp thụ độ chênh lệch số lượng tử giữa hai trạng thái nội tại, nó có thể trộn hai trạng thái ấy. Dao động tạo thành có thể biểu hiện dưới dạng khối lượng.

\section{Vì sao Higgs không giải thích mọi khối lượng}

Trước khi áp dụng khuôn mẫu đó, cần gạt bỏ một quan niệm sai lầm. Nói rằng mọi vật thể có khối lượng đều nhận toàn bộ khối lượng từ trường Higgs là không đúng.

Cho bức xạ vào một chiếc hộp phản xạ hoàn hảo. Từng photon đều không có khối lượng, nhưng chiếc hộp kín có hệ quy chiếu nghỉ và chứa năng lượng. Khối lượng của nó tăng thêm
\begin{equation}
\Delta M=\frac{E_{\mathrm{radiation}}}{c^{2}}.
\end{equation}
Vì vậy, khối lượng có thể bắt nguồn từ động năng và năng lượng trường bị giam giữ, ngay cả khi các thành phần cấu tạo không có khối lượng nghỉ.

Proton là ví dụ thực tế quan trọng. Nó chứa quark, phản quark và gluon trong một trạng thái tương tác mạnh. Phần lớn khối lượng proton đến từ động năng và năng lượng liên kết của chúng, chứ không phải từ phép cộng đơn giản khối lượng do Higgs tạo ra cho một vài quark cấu thành. Theo ước lượng được nhấn mạnh trong bài giảng, nếu làm cho các quark nhẹ không có khối lượng thì khối lượng proton chỉ thay đổi ở mức xấp xỉ một phần trăm. Lỗ đen làm nổi bật hơn nữa điểm cốt yếu về mặt logic: lỗ đen có thể có khối lượng nhờ tổng năng lượng đã tạo nên nó, chứ không phải vì trước đó mọi phần đóng góp đều được trường Higgs chuyển thành khối lượng nghỉ của các hạt cơ bản.

Cơ chế Higgs có một nhiệm vụ chính xác hơn. Nó cho phép fermion cơ bản và các boson chuẩn yếu có khối lượng, đồng thời vẫn giữ cho lý thuyết chuẩn nền tảng nhất quán.

\section{Tính chiral (tính thuận tay): khối lượng do sự trộn giữa trái--phải}

Trước hết, hãy tập trung vào electron. Spin là mômen động lượng nội tại. Với một hạt đang chuyển động, ta có thể so sánh hướng spin với động lượng của nó. Theo cách diễn giải hình học trong bài giảng, electron thuận tay phải và electron thuận tay trái được phân biệt qua mối liên hệ giữa hai hướng ấy. Với hạt không khối lượng, helicity và tính chiral trùng khớp theo cách đơn giản cần cho lập luận; với hạt có khối lượng, cần phân biệt chính xác hai khái niệm này.

Hãy hình dung ta liên tục vượt qua một electron có khối lượng. Trong một hệ quy chiếu quán tính khác, động lượng của nó đổi hướng còn spin không đổi, nên tính thuận tay gắn với hướng chuyển động có thể đảo ngược. Không người quan sát nào có thể vượt qua một hạt hoàn toàn không có khối lượng, vì nó luôn chuyển động với tốc độ ánh sáng. Từ đó, bài giảng đưa ra nhận định mang tính trực giác: electron không khối lượng giữ nguyên một trạng thái chiral, còn electron có khối lượng liên tục trộn lẫn trạng thái trái và phải.

Phát biểu chính xác trong lý thuyết trường là số hạng khối lượng Dirac liên kết hai trường chiral,
\begin{equation}
\mathcal L_m
=
-m_e\bar e e
=
-m_e
\left(
  \bar e_L e_R
  +
  \bar e_R e_L
\right).
\end{equation}
\editorialnote{Số hạng khối lượng Dirac là cách hoàn chỉnh toán học tiêu chuẩn cho hình ảnh ``tốc độ lật'' trong bài giảng. Mức độ trộn chiral tỉ lệ với khối lượng; tốc độ theo thời gian theo nghĩa đen còn phụ thuộc vào trạng thái và sự giãn nở thời gian tương đối tính.}

Nếu $m_e=0$, các trường trái và phải truyền đi độc lập. Nếu $m_e\ne0$, các phương trình liên kết chúng. Theo nghĩa chính xác đó, khối lượng đo cường độ trộn giữa trái--phải.

\section{Định luật bảo toàn cấm số hạng khối lượng trần}

Sao không viết số hạng trộn ngay từ đầu? Tương tác yếu đối xử khác nhau với $e_L$ và $e_R$. Susskind đặt cho sự khác biệt liên quan một cái tên cố ý đơn giản: \emph{zilch}. Theo quy ước chuẩn hóa dùng trong bài giảng,
\begin{equation}
z(e_L)=1,
\qquad
z(e_R)=0.
\end{equation}
Hai trạng thái chiral có cùng điện tích nhưng khác số lượng tử yếu. Vì vậy, chuyển trạng thái trực tiếp
\begin{equation}
e_L\longrightarrow e_R
\end{equation}
sẽ không bảo toàn đại lượng do tương tác yếu mang. Trong lý thuyết chưa phá vỡ đối xứng, số hạng khối lượng trần bị cấm.

\editorialnote{``Zilch'' là cách nói tắt nhằm phục vụ giảng dạy. Trong Mô hình Chuẩn, các fermion trái và phải thuộc những biểu diễn khác nhau của $SU(2)_L\times U(1)_Y$, nên số hạng khối lượng Dirac trực tiếp không bất biến chuẩn. Các giá trị được gán trong bài giảng, $1$ và $0$, mã hóa sự khác biệt này mà không cần đưa ra đầy đủ bảng điện tích điện yếu.}

Bảng trình bày trở ngại này dưới dạng một chuỗi quá trình. Electron phải không thể chuyển thành electron trái nếu không có thứ gì khác mang đi phần chênh lệch zilch. Trở ngại tương tự cũng áp dụng cho quark, muon và tau; khi trường còn thiếu được bổ sung, cường độ tương tác sẽ khác nhau tùy loài hạt.

\begin{figure}[tbp]
\centering
\includegraphics[width=0.95\textwidth]{lecture_144_figure_09.png}
\caption{Chuỗi chiral trên bảng: một lượng tử của trạng thái ngưng tụ được đưa vào tại dấu gạch chéo màu xanh lá, cho phép nhãn zilch thay đổi giữa quá trình truyền của trạng thái phải và trạng thái trái.}
\lectureframe{00:50:15}
\end{figure}

\section{Chất ngưng tụ ``Ziggs'' và khối lượng fermion}

Đưa vào một trường có các lượng tử mang đúng số lượng tử yếu còn thiếu. Để tránh nhầm lẫn quá sớm với hạt Higgs vật lý, bài giảng gọi lượng tử liên quan của trường này là \emph{Ziggs}. Giả sử trường Ziggs có thế năng hình mũ Mexico và chân không của nó là một trạng thái ngưng tụ. Khi đó, electron trái có thể phát ra một lượng tử Ziggs rồi chuyển thành electron phải,
\begin{equation}
e_L\longrightarrow e_R+\text{Ziggs},
\end{equation}
hoặc electron phải có thể hấp thụ một lượng tử Ziggs rồi chuyển thành electron trái,
\begin{equation}
e_R+\text{Ziggs}\longrightarrow e_L.
\end{equation}
Trạng thái ngưng tụ hấp thụ hoặc cung cấp phần chênh lệch số lượng tử mà không làm thay đổi trạng thái vĩ mô của nó. Điều bị cấm trong chân không trống rỗng chưa phá vỡ đối xứng trở thành khả dĩ trong chân không có đối xứng bị phá vỡ.

Sự phát xạ và hấp thụ lặp đi lặp lại làm trộn $e_L$ và $e_R$. Vì vậy, electron truyền đi như một hạt có khối lượng. Cơ chế tương tự cũng áp dụng cho các lepton mang điện và quark. Mỗi loài hạt có mức liên kết riêng với trạng thái ngưng tụ, nên cường độ chuyển trạng thái, và do đó khối lượng, có thể khác nhau giữa các loài.

Theo ký hiệu hiện đại, tương tác bất biến chuẩn này là tương tác Yukawa. Sau khi trường Higgs có giá trị kỳ vọng chân không $v$, tương tác tạo ra
\begin{equation}
\mathcal L_Y
\supset
-\frac{y_f}{\sqrt{2}}
(v+h)\bar f f,
\end{equation}
và do đó
\begin{equation}
m_f=\frac{y_fv}{\sqrt{2}},
\qquad
g_{hff}=\frac{y_f}{\sqrt{2}}
=
\frac{m_f}{v}.
\end{equation}
\editorialnote{Các công thức này đưa vào ký hiệu điện yếu tiêu chuẩn cho lập luận về sự liên kết với trạng thái ngưng tụ trong bài giảng. Chúng cũng làm rõ vì sao cùng một hằng số liên kết chi phối cả khối lượng fermion lẫn sự phân rã của Higgs thành fermion ấy.}

Bài giảng mô tả đây là sự phá vỡ tự phát đối xứng chiral. Mô tả đó nắm bắt được câu chuyện trái--phải trong phạm vi ví dụ đang xét. Chính xác hơn, giá trị kỳ vọng của Higgs phá vỡ đối xứng chuẩn điện yếu xuống nhóm con điện từ, rồi các tương tác Yukawa trở thành những số hạng khối lượng fermion.

\section{Boson yếu có khối lượng như thế nào}

Trạng thái ngưng tụ cũng làm thay đổi các boson chuẩn yếu. Một kích thích Ziggs có thể phát ra boson $Z$; theo chiều ngược lại, boson $Z$ có thể chuyển thành kích thích thích hợp được trạng thái ngưng tụ hỗ trợ rồi chuyển trở lại. Vì thế, boson $Z$ truyền đi trong chân không nơi nó trộn với khu vực góc của trường Higgs. Kết quả là một boson vectơ có khối lượng và có phân cực dọc mà boson vectơ không khối lượng không có.

\begin{figure}[tbp]
\centering
\includegraphics[width=0.95\textwidth]{lecture_144_figure_10.png}
\caption{Chuỗi hoàn chỉnh trên bảng về cơ chế của boson yếu: chuỗi chiral vẫn nằm phía trên, còn boson $Z$ trao đổi số lượng tử với khu vực phía dưới được trạng thái ngưng tụ hỗ trợ.}
\lectureframe{00:53:40}
\end{figure}

Đây là hiện tượng Brout--Englert--Higgs ở dạng dành cho boson chuẩn. Trong Mô hình Chuẩn đầy đủ, cả boson $W$ lẫn $Z$ đều có khối lượng, trong khi một tổ hợp của các trường chuẩn trung hòa ban đầu vẫn là photon không khối lượng. Nếu chân không phá vỡ đối xứng chuẩn điện từ thông thường, photon sẽ có khối lượng, như thực tế nó có bên trong chất siêu dẫn. Khi đó, tính chất điện từ tầm xa của không gian trống sẽ hoàn toàn khác.

Các boson yếu đã cung cấp bằng chứng mạnh mẽ cho cơ chế này từ lâu trước khi hạt Higgs vật lý được tách riêng. Khối lượng và tương tác của chúng cho thấy đối xứng điện yếu đã bị phá vỡ theo đúng cách cần thiết. Vì thế, theo nghĩa thực nghiệm không quá chặt chẽ, bài giảng có thể nói rằng khu vực ``Ziggs'' đã được biết đến: đối tượng thực nghiệm còn thiếu là dao động bán kính độc lập.

\editorialnote{Lời nhắc ứng khẩu trong bài giảng về một phát hiện tại SLAC vào khoảng 1980 gộp làm một nhiều sự kiện. Dòng trung hòa yếu được quan sát vào năm 1973, tán xạ electron vi phạm tính chẵn lẻ tại SLAC được báo cáo vào cuối thập niên 1970, còn các boson $W$ và $Z$ được phát hiện trực tiếp tại CERN vào năm 1983. Điểm vật lý cốt yếu là sự phá vỡ đối xứng điện yếu đã được thực nghiệm xác nhận từ trước.}

\section{Bản chất thực sự của boson Higgs}

Hãy quay lại với chiếc mũ Mexico. Chuyển động quanh rãnh làm thay đổi góc nhưng không đổi bán kính. Khu vực góc này cung cấp các mode được hấp thụ vào các boson yếu có khối lượng. Boson Higgs vật lý tương ứng với kiểu chuyển động kia: dao động vào trong rồi ra ngoài, đi lên các thành của thế năng.

Có thể viết trường dưới dạng sơ đồ như sau
\begin{equation}
\Phi(x)
=
\frac{v+h(x)}{\sqrt{2}}
\exp\!\left(
  \frac{i\theta(x)}{v}
\right).
\end{equation}
Biến $\theta$ mô tả chuyển động dọc theo rãnh. Biến $h$ làm thay đổi bán kính. Trong lý thuyết chuẩn, các bậc tự do góc thích hợp trở thành các thành phần dọc của boson yếu có khối lượng, còn $h$ vẫn là một hạt vô hướng truyền đi được: boson Higgs.

\begin{figure}[tbp]
\centering
\includegraphics[width=0.95\textwidth]{lecture_144_figure_11.png}
\caption{Cả hai thế năng lại xuất hiện trên bảng, với độ dịch chuyển theo phương bán kính được đánh dấu trên rãnh của thế năng hình mũ Mexico. Dao động vào trong và ra ngoài này là mode Higgs vật lý.}
\lectureframe{00:56:30}
\end{figure}

Có một cách hình dung tương đương dùng nhiều hạt. Chân không là trạng thái ngưng tụ, và dao động theo phương bán kính làm thay đổi mật độ của trạng thái ấy. Do đó, boson Higgs giống như một sóng nén trong môi trường lượng tử đó. Phép tương tự này hữu ích miễn là ta không suy ra rằng sóng nhất thiết phải bền. Sóng âm thông thường có thể phân rã, và kích thích Higgs cũng vậy.

Sự phân biệt này giải thích vì sao việc tìm thấy boson Higgs là mảnh ghép cuối cùng còn thiếu, chứ không phải bằng chứng đầu tiên cho cơ chế. Chân không bị phá vỡ đối xứng và các boson yếu có khối lượng đã được nhận thấy gián tiếp từ trước. Cần tạo ra sóng mật độ theo phương bán kính dưới dạng hạt rồi nhận diện nó qua các sản phẩm phân rã.

\section{Phân rã, tạo sinh và lợi thế của fermion nặng}

Higgs liên kết với một fermion với cường độ tỉ lệ với khối lượng của fermion đó. Nó có thể phân rã thành một cặp fermion-phản fermion,
\begin{equation}
H\longrightarrow f+\bar f,
\end{equation}
và cũng có thể đọc cùng sơ đồ ấy theo chiều ngược lại như quá trình tạo sinh,
\begin{equation}
f+\bar f\longrightarrow H.
\end{equation}
Vì vậy, về nguyên tắc, hủy cặp electron-positron có thể tạo ra Higgs. Tuy nhiên, trong lịch sử, các máy gia tốc đời đầu chưa đủ năng lượng để tạo ra Higgs; về sau, các máy gặp một trở ngại cơ bản hơn: tương tác Yukawa của electron cực nhỏ. Quá trình hủy cặp quark nhẹ cũng yếu tương tự.

\begin{figure}[tbp]
\centering
\includegraphics[width=0.95\textwidth]{lecture_144_figure_12.png}
\caption{Đỉnh phân rã Higgs được vẽ trên bảng. Đọc các mũi tên theo chiều ngược lại ta được sự hợp nhất của fermion và phản fermion thành Higgs.}
\lectureframe{00:59:40}
\end{figure}

Fermion càng nặng thì liên kết Yukawa của nó càng mạnh. Điều này làm nổi bật quark top, có khối lượng khoảng $173\,\mathrm{GeV}$, tức xấp xỉ $180$ lần khối lượng proton. Bài giảng dùng con số tròn khoảng $170$ lần khối lượng proton. Tuy vậy, Higgs có khối lượng gần $125\,\mathrm{GeV}$ thì quá nhẹ để phân rã thành một cặp $t\bar t$ trên vỏ khối lượng,
\begin{equation}
m_H<2m_t.
\end{equation}
Cũng không thể tạo ra một chùm quark top hữu ích: quark top phân rã trước khi kịp hadron hóa thành hạt có tuổi thọ đủ dài để tạo thành chùm.

Cơ học lượng tử cung cấp cách vượt qua cả hai trở ngại. Hai gluon có thể dao động qua một vòng ảo của quark top rồi xuất hiện dưới dạng một Higgs,
\begin{equation}
g+g\longrightarrow H.
\end{equation}
Các đường quark top ảo không cần thỏa ngưỡng trên vỏ khối lượng vốn chi phối quá trình tạo cặp quark top thực. Proton chứa đầy gluon, vì vậy va chạm proton-proton năng lượng cao tại Máy Va chạm Hadron Lớn cung cấp nguồn dồi dào các cặp gluon cần thiết.

\begin{figure}[tbp]
\centering
\includegraphics[width=0.95\textwidth]{lecture_144_figure_13.png}
\caption{Sự hợp nhất gluon qua một vòng quark top ảo. Hai gluon lấy từ các proton va chạm gặp nhau trong vòng và tạo ra Higgs.}
\lectureframe{01:06:35}
\end{figure}

Đó là lý do máy va chạm proton, dù tạo ra đám hạt thứ cấp phức tạp, đã trở thành cơ sở sản xuất Higgs mang tính quyết định. Proton không tương tác với Higgs như một vật thể sơ cấp duy nhất; các thành phần gluon của nó tạo ra Higgs qua một vòng, trong đó liên kết mạnh với quark top nặng chi phối quá trình.

\section{Hạt 2012 và phép thử hai photon}

Hạt được công bố vào năm 2012 có khối lượng gần
\begin{equation}
m_H\simeq125\,\mathrm{GeV}.
\end{equation}
Theo phép so sánh thuận tiện trong bài giảng, con số này xấp xỉ $127$ lần khối lượng proton. Việc phát hiện hạt này hoàn tất danh sách các hạt mà Mô hình Chuẩn tối thiểu đòi hỏi. Câu hỏi còn lại là liệu các tính chất định lượng của nó có đúng như những gì Mô hình Chuẩn dự đoán hay không.

Có thể đảo ngược lập luận về vòng quark top. Higgs tương tác với các hạt mang điện ảo trong một vòng, và vòng đó có thể phát ra hai photon,
\begin{equation}
H\longrightarrow\gamma+\gamma.
\end{equation}
Phân rã này hiếm gặp, nhưng photon là những đối tượng đặc biệt sạch trong máy va chạm hadron. Vì vậy, một đỉnh dư hẹp trong phân bố khối lượng bất biến của chúng có thể làm lộ ra Higgs, ngay cả khi các phân rã phổ biến hơn khó phân tích thực nghiệm hơn.

\begin{figure}[tbp]
\centering
\includegraphics[width=0.95\textwidth]{lecture_144_figure_14.png}
\caption{Tam giác hai photon ở bên phải bảng, đặt cạnh vòng hợp nhất gluon đã trình bày trước đó ở bên trái. Hai quá trình khảo sát cùng một lập luận ở cấp độ vòng theo hai chiều ngược nhau.}
\lectureframe{01:09:00}
\end{figure}

Vào thời điểm diễn ra bài giảng, tốc độ phân rã thành hai photon được báo cáo dường như cao gấp khoảng một phẩy năm lần so với dự đoán đơn giản nhất của Mô hình Chuẩn, với độ ý nghĩa chỉ khoảng hai độ lệch chuẩn. Chỉ riêng dao động như vậy chưa phải là bằng chứng. Tuy nhiên, nếu mức này vẫn tồn tại, nó có thể chỉ ra một hạt mang điện bổ sung trong vòng, có lẽ liên quan đến siêu đối xứng hoặc một mở rộng khác của Mô hình Chuẩn.

Nhận định này thuộc về bối cảnh lịch sử của năm 2012. Nó ghi lại điều mà các nhà thực nghiệm được dặn là hãy theo dõi tiếp, chứ không phải là khẳng định về một dị thường hiện nay. Bản thân phát hiện này phần lớn xác nhận một khuôn khổ vốn đã được dữ liệu về các boson yếu ủng hộ. Những phép thử khoa học sắc bén hơn nằm ở các tốc độ phân rã: liệu hạt vô hướng mới có tương tác với từng hạt theo quy luật phụ thuộc khối lượng mà cơ chế Higgs dự đoán hay không?

\section{Câu hỏi từ thính giả}

\begin{classroomqa}
\audiencequestion{Điều gì quyết định vì sao fermion này chuyển trạng thái giữa trái và phải nhanh hơn fermion khác, và do đó có khối lượng khác?\lecturetimestamp{01:12:00}}
\lecturerresponse{Mỗi loại fermion có một liên kết riêng với chất ngưng tụ Higgs. Liên kết mạnh hơn tạo ra sự trộn trái--phải mạnh hơn và khối lượng lớn hơn. Cũng chính liên kết đó quyết định mức độ mạnh yếu của quá trình Higgs vật lý phân rã thành fermion ấy; vì vậy, fermion nặng hơn nhìn chung tương tác mạnh hơn. Mô hình Chuẩn mô tả các liên kết Yukawa này bằng các tham số nhưng không giải thích được thứ bậc phân cấp nổi bật của chúng. Nói rằng electron có liên kết yếu còn quark top có liên kết mạnh là mô tả chính xác, nhưng chưa phải là một giải thích sâu hơn.}
\end{classroomqa}

\begin{classroomqa}
\audiencequestion{Ta có biết độ lớn của trường Higgs trong chân không không?\lecturetimestamp{01:13:36}}
\lecturerresponse{Có. Theo quy ước dùng ở đây, giá trị kỳ vọng chân không vào khoảng $240\,\mathrm{GeV}$; giá trị chuẩn chính xác hơn xấp xỉ $246\,\mathrm{GeV}$. Có thể hình dung nó như độ dịch chuyển theo phương bán kính của trường so với gốc tọa độ, hoặc như mật độ của chất ngưng tụ. Một lượng tử Higgs vật lý là dao động nhỏ theo phương bán kính quanh giá trị ấy, tương đương với một biến thiên cục bộ về mật độ của chất ngưng tụ.}
\end{classroomqa}

\begin{classroomqa}
\audiencequestion{Giá trị chân không khác không của trường Higgs có hàm ý rằng không gian trống có mật độ năng lượng không?\lecturetimestamp{01:14:12}}
\lecturerresponse{Có. Thế năng Higgs đóng góp vào năng lượng chân không, cũng như nhiều sector khác của lý thuyết trường lượng tử. Năng lượng chân không tổng thể được quan sát nhỏ khác thường so với độ lớn tự nhiên của các đóng góp ấy, nên chúng phải gần như triệt tiêu lẫn nhau. Chưa có lý thuyết được chấp nhận nào giải thích sự triệt tiêu đó. Đây là vấn đề hằng số vũ trụ, một trong những vấn đề sâu sắc nhất chưa được giải quyết trong vật lý cơ bản.}
\end{classroomqa}

\section{Cấu trúc này đã đạt được điều gì}

Giờ đây có thể lần theo lập luận từ đầu đến cuối mà không cần viện đến một thứ ê-te dính nhớp. Phép quay nội tại lượng tử hóa cung cấp mô hình cho điện tích. Biến cái bát thành một chiếc mũ Mexico khiến trường có giá trị chân không khác không. Nguyên lý bất định biến chân không cổ điển được chọn sắc nét thành một chất ngưng tụ không nhạy với từng lượng tử riêng lẻ. Khi đó, chất ngưng tụ có thể mang độ chênh số lượng tử yếu giữa các fermion tay trái và tay phải, cho phép chúng trộn lẫn và do đó có khối lượng. Nó cũng cung cấp các mode dọc, nhờ đó $W$ và $Z$ truyền đi như những boson vectơ có khối lượng, trong khi photon vẫn không có khối lượng.

Boson Higgs vật lý là dao động theo phương bán kính còn lại sau sự tái sắp xếp này. Nó không phải là ``bản thân khối lượng'', và cũng không phải là nguyên nhân tạo ra phần lớn khối lượng của proton. Nó là sóng mật độ quan sát được của trường phá vỡ đối xứng. Các liên kết của nó lưu giữ dấu ấn của những khối lượng do trường đó tạo ra; vì thế, các hạt nặng chi phối quá trình sinh Higgs và phân rã của nó. Do đó, vòng hợp nhất gluon, kênh hai photon và câu hỏi cuối của thính giả về năng lượng chân không không phải là những điều hiếu kỳ tách rời nhau. Chúng là những hệ quả, phép thử và hàm ý còn bỏ ngỏ của cùng một chân không phi tầm thường.
